\documentclass[10pt,a4paper]{amsart}
\usepackage{amsmath,amssymb,amsthm,graphicx}
\usepackage{verbatim}
\newenvironment{reference}{\comment}{\endcomment}
%\newenvironment{reference}{}{}
\newenvironment{slogan}{\comment}{\endcomment}
\newenvironment{history}{\comment}{\endcomment}

% For commutative diagrams we use Xy-pic
\usepackage[all]{xy}

% We use 2cell for 2-commutative diagrams.
\xyoption{2cell}
% Theorem environments.
%
\theoremstyle{plain}
\newtheorem{theorem}[subsection]{Theorem}
\newtheorem{proposition}[subsection]{Proposition}
\newtheorem{lemma}[subsection]{Lemma}

\theoremstyle{definition}
\newtheorem{definition}[subsection]{Definition}
\newtheorem{example}[subsection]{Example}
\newtheorem{exercise}[subsection]{Exercise}
\newtheorem{situation}[subsection]{Situation}

\theoremstyle{remark}
\newtheorem{remark}[subsection]{Remark}
\newtheorem{remarks}[subsection]{Remarks}

\numberwithin{equation}{subsection}

% Macros
%
\def\lim{\mathop{\mathrm{lim}}\nolimits}
\def\colim{\mathop{\mathrm{colim}}\nolimits}
\def\Spec{\mathop{\mathrm{Spec}}}
\def\Hom{\mathop{\mathrm{Hom}}\nolimits}
\def\Ext{\mathop{\mathrm{Ext}}\nolimits}
\def\SheafHom{\mathop{\mathcal{H}\!\mathit{om}}\nolimits}
\def\SheafExt{\mathop{\mathcal{E}\!\mathit{xt}}\nolimits}
\def\Sch{\mathit{Sch}}
\def\Mor{\mathop{\mathrm{Mor}}\nolimits}
\def\Ob{\mathop{\mathrm{Ob}}\nolimits}
\def\Sh{\mathop{\mathit{Sh}}\nolimits}
\def\NL{\mathop{N\!L}\nolimits}
\def\CH{\mathop{\mathrm{CH}}\nolimits}
\def\proetale{{pro\text{-}\acute{e}tale}}
\def\etale{{\acute{e}tale}}
\def\QCoh{\mathit{QCoh}}
\def\Ker{\mathop{\mathrm{Ker}}}
\def\Im{\mathop{\mathrm{Im}}}
\def\Coker{\mathop{\mathrm{Coker}}}
\def\Coim{\mathop{\mathrm{Coim}}}

% Boxtimes
%
\DeclareMathSymbol{\boxtimes}{\mathbin}{AMSa}{"02}

%
% Macros for moduli stacks/spaces
%
\def\QCohstack{\mathcal{QC}\!\mathit{oh}}
\def\Cohstack{\mathcal{C}\!\mathit{oh}}
\def\Spacesstack{\mathcal{S}\!\mathit{paces}}
\def\Quotfunctor{\mathrm{Quot}}
\def\Hilbfunctor{\mathrm{Hilb}}
\def\Curvesstack{\mathcal{C}\!\mathit{urves}}
\def\Polarizedstack{\mathcal{P}\!\mathit{olarized}}
\def\Complexesstack{\mathcal{C}\!\mathit{omplexes}}
% \Pic is the operator that assigns to X its picard group, usage \Pic(X)
% \Picardstack_{X/B} denotes the Picard stack of X over B
% \Picardfunctor_{X/B} denotes the Picard functor of X over B
\def\Pic{\mathop{\mathrm{Pic}}\nolimits}
\def\Picardstack{\mathcal{P}\!\mathit{ic}}
\def\Picardfunctor{\mathrm{Pic}}
\def\Deformationcategory{\mathcal{D}\!\mathit{ef}}
\usepackage{cleveref}
\crefname{originallemma}{original lemma}{original lemmas}
\Crefname{originallemma}{Original lemma}{Original lemmas}
\makeatletter
\newlabel{lemma-maximal-subfield-splits}{{0750}{source}{}{sourcepointer-1}{}}
\newlabel{lemma-maximal-subfield-splits@cref}{{[originallemma][0][]0750}{[1][0][]source}{}{}{}}
\makeatother
\crefname{originalsection}{original section}{original sections}
\Crefname{originalsection}{Original section}{Original sections}
\makeatletter
\newlabel{fields-section-algebraic}{{037H}{source}{}{sourcepointer-2}{}}
\newlabel{fields-section-algebraic@cref}{{[originalsection][0][]037H}{[1][0][]source}{}{}{}}
\makeatother
\crefname{originallemma}{original lemma}{original lemmas}
\Crefname{originallemma}{Original lemma}{Original lemmas}
\makeatletter
\newlabel{lemma-base-change}{{074H}{source}{}{sourcepointer-3}{}}
\newlabel{lemma-base-change@cref}{{[originallemma][0][]074H}{[1][0][]source}{}{}{}}
\makeatother
\crefname{originallemma}{original lemma}{original lemmas}
\Crefname{originallemma}{Original lemma}{Original lemmas}
\makeatletter
\newlabel{lemma-brauer-algebraically-closed}{{074M}{source}{}{sourcepointer-4}{}}
\newlabel{lemma-brauer-algebraically-closed@cref}{{[originallemma][0][]074M}{[1][0][]source}{}{}{}}
\makeatother
\title{A finite central skew field has a maximal separable subfield}
\author{The Stacks Project Authors}
\date{stacks-project:brauer.tex:local-snapshot-20261012; GNU FDL 1.2 or later}
\begin{document}
\maketitle
\noindent\textit{Editorial note.} All delivered mathematical text is the author's own, copied from the pinned Stacks Project chapter file sites-cohomology.tex; no statement or proof step was rewritten, reordered or removed.\par\medskip
\noindent\textit{Editorial note.} The proof is the author's complete argument; the single parenthetical '(details omitted)' concerns the routine claim that the coordinate functions of the q-th power map are polynomials, which follows from bilinearity of the multiplication over the ground field and is stated in the counted span immediately before it, all the other steps being written out. The parenthetical is preserved verbatim.\par\medskip
\noindent\textit{Editorial note.} Every labelled reference occurring inside the counted spans is carried in assembly.external\_references as a non-load-bearing navigation pointer (same-chapter ones at their definition site, cross-chapter ones at the invocation inside the counted span), so no load-bearing mathematics is replaced by a stub. No bibliography citation occurs inside the counted spans.\par\medskip
\noindent\textit{Editorial note.} The same-name check for this card scanned every staged Stacks chapter of this repository -- all sources/stacks-*/pinned-original-source.tex files, the card's own chapter excluded -- and the labels found are recorded in dedup.same\_name\_cross\_chapter on the card.\par\medskip
\noindent\textit{Editorial note.} The preamble is the chapter's own retained minimal preamble (sources/stacks-sites-cohomology/retained-preamble.tex), a verbatim excerpt of the native preamble: the Xy-pic setup and the native theorem declarations with the mathematical macros.\par\medskip
\noindent\textit{Allowed standard prerequisites (editorial).}\begin{itemize}
\item A maximal subfield splits the algebra. Same chapter, Lemma lemma-maximal-subfield-splits. Supplies that a finite central skew field is split by a maximal subfield, which turns the second assertion of the statement into the first.
\item The centralizer of a separable subfield. Same chapter, the discussion of centralizers of subfields of a central skew field. Supplies that the centralizer of a separable subfield again has that subfield as its center, which is what allows the author to replace the ground field and reduce to finding one separable element.
\item Inseparable elements and minimal polynomials. Fields, Section fields-section-algebraic. Supplies that in positive characteristic an element whose minimal polynomial is not separable has a minimal polynomial of the form a p-power variable minus a constant.
\item Base change and the algebraic closure case. Same chapter, Lemma lemma-base-change and Lemma lemma-brauer-algebraically-closed. Supply that extending scalars to the algebraic closure turns the skew field into a matrix algebra, which produces the element whose p-th power is not central.
\end{itemize}\medskip
\ifdefined\hypertarget\hypertarget{sourcepointer-1}{}\fi
\noindent\textit{Source reference.} Original lemma 0750: sources/stacks-brauer/pinned-original-source.tex, source lines 896--896. Navigation pointer invoked at proof line 935; the counted span contains the pointer only and the cited result is not reproduced inside it, so the pointer is registered as a non-load-bearing navigation stub in the established form for Stacks material. No bibliography citation is involved.\par\medskip
\ifdefined\hypertarget\hypertarget{sourcepointer-2}{}\fi
\noindent\textit{Source reference.} Original section 037H: sources/stacks-fields/pinned-original-source.tex, source lines 3909--3909. Navigation pointer invoked at proof line 952; the counted span contains the pointer only and the cited result is not reproduced inside it, so the pointer is registered as a non-load-bearing navigation stub in the established form for Stacks material. No bibliography citation is involved.\par\medskip
\ifdefined\hypertarget\hypertarget{sourcepointer-3}{}\fi
\noindent\textit{Source reference.} Original lemma 074H: sources/stacks-brauer/pinned-original-source.tex, source lines 565--565. Navigation pointer invoked at proof line 975; the counted span contains the pointer only and the cited result is not reproduced inside it, so the pointer is registered as a non-load-bearing navigation stub in the established form for Stacks material. No bibliography citation is involved.\par\medskip
\ifdefined\hypertarget\hypertarget{sourcepointer-4}{}\fi
\noindent\textit{Source reference.} Original lemma 074M: sources/stacks-brauer/pinned-original-source.tex, source lines 667--667. Navigation pointer invoked at proof line 975; the counted span contains the pointer only and the cited result is not reproduced inside it, so the pointer is registered as a non-load-bearing navigation stub in the established form for Stacks material. No bibliography citation is involved.\par\medskip
\begin{proposition}
\label{proposition-separable-splitting-field}
Consider a finite central skew field $K$ over $k$.
There exists a maximal subfield $k \subset k' \subset K$ which
is separable over $k$.
In particular, every Brauer class has a finite separable
spitting field.
\end{proposition}

\begin{proof}
Since every Brauer class is represented by a finite central skew
field over $k$, we see that the second statement follows from the
first by
Lemma \ref{lemma-maximal-subfield-splits}.

\medskip\noindent
To prove the first statement, suppose that we are given a separable
subfield $k' \subset K$. Then the centralizer $K'$ of $k'$ in $K$
has center $k'$, and the problem reduces to finding a maximal
subfield of $K'$ separable over $k'$. Thus it suffices to prove, if
$k \not = K$, that we can find an element $x \in K$, $x \not \in k$
which is separable over $k$. This statement is clear in characteristic
zero. Hence we may assume that $k$ has characteristic $p > 0$. If the
ground field $k$ is finite then, the result is clear as well (because
extensions of finite fields are always separable). Thus we may assume
that $k$ is an infinite field of positive characteristic.

\medskip\noindent
To get a contradiction assume no element of $K$ is separable over $k$.
By the discussion in
Fields, Section \ref{fields-section-algebraic}
this means the minimal polynomial of any $x \in K$ is of the form
$T^q - a$ where $q$ is a power of $p$ and $a \in k$. Since it is
clear that every element of $K$ has a minimal polynomial of degree
$\leq \dim_k(K)$ we conclude that there exists a fixed $p$-power
$q$ such that $x^q \in k$ for all $x \in K$.

\medskip\noindent
Consider the map
$$
(-)^q : K \longrightarrow K
$$
and write it out in terms of a $k$-basis $\{a_1, \ldots, a_n\}$ of $K$
with $a_1 = 1$. So
$$
(\sum x_i a_i)^q = \sum f_i(x_1, \ldots, x_n)a_i.
$$
Since multiplication on $K$ is $k$-bilinear we see that each $f_i$
is a polynomial in $x_1, \ldots, x_n$ (details omitted).
The choice of $q$ above and the fact that $k$ is infinite shows that
$f_i$ is identically zero for $i \geq 2$. Hence we see that it remains
zero on extending $k$ to its algebraic closure $\overline{k}$. But the
algebra $K \otimes_k \overline{k}$ is a matrix algebra (for example by
Lemmas \ref{lemma-base-change} and \ref{lemma-brauer-algebraically-closed}),
which implies there are some elements whose $q$th power is not central
(e.g., $e_{11}$). This is the desired contradiction.
\end{proof}

\end{document}
